\documentclass[12pt]{article}
\usepackage[a4paper,margin=1in]{geometry}
\usepackage[utf8]{inputenc}
\usepackage[T1]{fontenc}
\usepackage{amsmath,amssymb}
\usepackage{textcomp}
\usepackage{parskip}
\begin{document}
\begin{center}
{\Large \textbf{2013 STEP III}}
\end{center}

\subsection*{Page 1}
\begin{flushleft}
\bigskip\textbf{Section A: Pure Mathematics}\par
dt 2t\par
\bigskip\textbf{Question 1.} Given that t = tan 1 x, show that = 1 (1 + t2) and sin x = .\par
\[2 dx 2 1 + t^2\]
Hence show that\par
\[2 \pi 1 2 \sqrt\{1\} a\]
\[dx = arctan - (0 < a < 1).\]
\[1 + a sin x \sqrt\{1\} a^2 \sqrt\{1\} + a\]
\[\leq0\]
-\par
Let\par
\[2 \pi sinn x\]
\[I = dx (n \leq 0).\]
n\par
2 + sin x\par
\[\leq0\]
By considering I + 2I , or otherwise, evaluate I .\par
n+1 n 3\par
\bigskip\textbf{Question 2.} In this question, you may ignore questions of convergence.\par
arcsin x\par
\[Let y = . Show that\]
\[\sqrt\{1\} x^2\]
- dy\par
\[(1 x^2) xy 1 = 0\]
- dx - -\par
and prove that, for any positive integer n,\par
dn+2y dn+1y dny\par
\[(1 x^2) (2n + 3)x (n + 1)2 = 0 .\]
- dxn+2 - dxn+1 - dxn\par
arcsin x\par
Hence obtain the Maclaurin series for , giving the general term for odd and for even\par
\[\sqrt\{1\} x^2\]
powers of x. -\par
Evaluate the infinite sum\par
\[1 22 22 32 22 32 n^2\]
\[1 + + + \times + + \times \times \cdots \times + .\]
3! 5! 7! \cdots (2n + 1)! \cdots\par
\end{flushleft}
\subsection*{Page 2}
\begin{flushleft}
\bigskip\textbf{Question 3.} The four vertices P (i = 1, 2, 3, 4) of a regular tetrahedron lie on the surface of a sphere with\par
i\par
centre at O and of radius 1. The position vector of P with respect to O is p (i = 1, 2, 3, 4).\par
i i\par
\[Use the fact that p + p + p + p = 0 to show that p . p = 1 for i = j.\]
\[1 2 3 4 i j - 3 \geq\]
Let X be any point on the surface of the sphere, and let XP denote the length of the line\par
i\par
\[joining X and P (i = 1, 2, 3, 4).\]
i\par
\textbf{(i)} By writing (XP )2 as (p x) . (p x), where x is the position vector of X with respect\par
i i i\par
- -\par
to O, show that\par
\[(XP )2 = 8 .\]
i\par
\[i=1\]
\[\geq\]
\textbf{(ii)} Given that P has coordinates (0, 0, 1) and that the coordinates of P are of the form\par
1 2\par
\[(a, 0, b), where a > 0, show that a = 2\sqrt\{2\}/3 and b = 1/3, and find the coordinates of\]
-\par
P and P .\par
3 4\par
\textbf{(iii)} Show that\par
4 4\par
\[(XP )4 = 4 (1 x . p )2 .\]
i i\par
-\par
\[i=1 i=1\]
\[\geq \geq\]
By letting the coordinates of X be (x, y, z), show further that (XP )4 is independent\par
i\par
\[i=1\]
of the position of X. (\par
\bigskip\textbf{Question 4.} Show that (z eiθ)(z e -iθ) = z2 2z cos θ + 1 .\par
- - -\par
\[Write down the (2n)th roots of 1 in the form eiθ, where \pi < θ \geq \pi, and deduce that\]
- -\par
n\par
\[z2n + 1 = z^2 2z cos (2k-1)\pi + 1 .\]
- 2n\par
\[k ) =1 \int \int \sum \sum\]
Here, n is a positive integer, and the notation denotes the product.\par
(\par
\textbf{(i)} By substituting z = i show that, when n is even,\par
\[cos \pi cos 3\pi cos 5\pi cos (2n-1)\pi = ( 1) 2 1 n 21-n .\]
2n 2n 2n \cdots 2n -\par
\[\int \sum\]
) [ ) [ ) [\par
\textbf{(ii)} Show that, when n is odd,\par
\[cos2 \pi cos2 3\pi cos2 5\pi cos2 (n-2)\pi = n^21-n .\]
2n 2n 2n \cdots 2n\par
\[\int \sum\]
) [ ) [ ) [\par
\[You may use without proof the fact that 1 + z2n = (1 + z^2)(1 z^2 + z^4 + z2n-2)\]
- - \cdots\par
when n is odd.\par
\end{flushleft}
\subsection*{Page 3}
\begin{flushleft}
\bigskip\textbf{Question 5.} In this question, you may assume that, if a, b and c are positive integers such that a and b\par
are coprime and a divides bc, then a divides c. (Two positive integers are said to be coprime\par
if their highest common factor is 1.)\par
\textbf{(i)} Suppose that there are positive integers p, q, n and N such that p and q are coprime\par
\[and qnN = pn. Show that N = kpn for some positive integer k and deduce the value\]
of q.\par
Hence prove that, for any positive integers n and N ,\par
\[\sqrt\{n\}\]
N is either a positive integer or\par
irrational.\par
\textbf{(ii)} Suppose that there are positive integers a, b, c and d such that a and b are coprime and\par
\[c and d are coprime, and aadb = bacb . Prove that db = ba and deduce that, if p is a\]
prime factor of d, then p is also a prime factor of b.\par
If pm and pn are the highest powers of the prime number p that divide d and b, respec-\par
\[tively, express b in terms of a, m and n and hence show that pn \geq n. Deduce the value\]
\[of b. (You may assume that if x > 0 and y \leq 2 then yx > x.)\]
Hence prove that, if r is a positive rational number such that rr is rational, then r is a\par
positive integer.\par
\bigskip\textbf{Question 6.} Let z and w be complex numbers. Use a diagram to show that z w \geq z + w .\par
| - | | | | |\par
For any complex numbers z and w, E is defined by\par
∗ ∗\par
\[E = zw + z w + 2 zw .\]
| |\par
\textbf{(i)} Show that z w 2 = ( z + w )2 E , and deduce that E is real and non-negative.\par
| - | | | | | -\par
\textbf{(ii)} Show that 1 zw ∗ 2 = (1 + zw )2 E .\par
| - | | | -\par
\[Hence show that, if both z > 1 and w > 1, then\]
| | | |\par
z w z + w\par
\[| - | \geq | | | | .\]
1 zw∗ 1 + zw\par
| - | | |\par
\[Does this inequality also hold if both z < 1 and w < 1?\]
| | | |\par
\end{flushleft}
\subsection*{Page 4}
\begin{flushleft}
\[d2y dy\]
\bigskip\textbf{Question 7.} (i) Let y(x) be a solution of the differential equation + y3 = 0 with y = 1 and = 0\par
\[dx2 dx\]
\[at x = 0, and let\]
dy\par
\[E(x) = + 1 y^4 .\]
dx 2\par
] \{\par
\[Show by differentiation that E is constant and deduce that y(x) \geq 1 for all x.\]
| |\par
\[d2v dv\]
\textbf{(ii)} Let v(x) be a solution of the differential equation + x + sinh v = 0 with v = ln 3\par
\[dx2 dx\]
dv\par
\[and = 0 at x = 0, and let\]
dx\par
dv\par
\[E(x) = + 2 cosh v .\]
dx\par
] \{\par
dE\par
\[Show that \geq 0 for x \leq 0 and deduce that cosh v(x) \geq 5 for x \leq 0.\]
dx 3\par
\textbf{(iii)} Let w(x) be a solution of the differential equation\par
\[d2w dw\]
\[+ (5 cosh x 4 sinh x 3) + (w cosh w + 2 sinh w) = 0\]
\[dx2 - - dx\]
dw 1\par
\[with = and w = 0 at x = 0. Show that cosh w(x) \geq 5 for x \leq 0.\]
\[dx \sqrt\{2\} 4\]
\end{flushleft}
\subsection*{Page 5}
\begin{flushleft}
n-1\par
\[8 Evaluate e2i(α+r\pi/n) where α is a fixed angle and n \leq 2.\]
\[r=0\]
\[\geq\]
The fixed point O is a distance d from a fixed line D. For any point P , let s be the distance\par
from P to D and let r be the distance from P to O. Write down an expression for s in terms\par
of d, r and the angle θ, where θ is as shown in the diagram below.\par
D\par
E\par
P s\par
r\par
θ\par
O d\par
\[The curve E shown in the diagram is such that, for any point P on E, the relation r = ks\]
\[holds, where k is a fixed number with 0 < k < 1.\]
Each of the n lines L , . . . , L passes through O and the angle between adjacent lines is \pi .\par
1 n n\par
The line L (j = 1, . . . , n) intersects E in two points forming a chord of length l . Show that,\par
j j\par
\[for n \leq 2,\]
\[n 1 (2 k^2)n\]
\[= - .\]
l 4kd\par
j\par
\[j=1\]
\[\geq\]
\end{flushleft}
\subsection*{Page 6}
\begin{flushleft}
\bigskip\textbf{Section B: Mechanics}\par
\bigskip\textbf{Question 9.} A sphere of radius R and uniform density ρ is floating in a large tank of liquid of uniform\par
s\par
density ρ. Given that the centre of the sphere is a distance x above the level of the liquid,\par
\[where x < R, show that the volume of liquid displaced is\]
\[\pi\]
\[(2R3 3R2x + x^3) .\]
3 -\par
The sphere is acted upon by two forces only: its weight and an upward force equal in magnitude\par
to the weight of the liquid it has displaced. Show that\par
\[4R3ρ (g + x¨) = (2R3 3R2x + x^3)ρg .\]
s\par
-\par
Given that the sphere is in equilibrium when x = 1 R, find ρ in terms of ρ. Find, in terms of\par
2 s\par
R and g, the period of small oscillations about this equilibrium position.\par
\bigskip\textbf{Question 10.} A uniform rod AB has mass M and length 2a. The point P lies on the rod a distance\par
a x from A. Show that the moment of inertia of the rod about an axis through P and\par
-\par
perpendicular to the rod is\par
\[1 M (a^2 + 3x2) .\]
The rod is free to rotate, in a horizontal plane, about a fixed vertical axis through P . Initially\par
the rod is at rest. The end B is struck by a particle of mass m moving horizontally with speed\par
u in a direction perpendicular to the rod. The coefficient of restitution between the rod and\par
the particle is e. Show that the angular velocity of the rod immediately after impact is\par
3mu(1 + e)(a + x)\par
.\par
\[M (a^2 + 3x2) + 3m(a + x)2\]
In the case m = 2M , find the value of x for which the angular velocity is greatest and show\par
that this angular velocity is u(1 + e)/a .\par
\end{flushleft}
\subsection*{Page 7}
\begin{flushleft}
\bigskip\textbf{Question 11.} An equilateral triangle, comprising three light rods each of length \sqrt3a, has a particle of mass m\par
attached to each of its vertices. The triangle is suspended horizontally from a point vertically\par
above its centre by three identical springs, so that the springs and rods form a tetrahedron.\par
Each spring has natural length a and modulus of elasticity kmg, and is light. Show that when\par
the springs make an angle θ with the horizontal the tension in each spring is\par
kmg(1 cos θ)\par
- .\par
cos θ\par
\[Given that the triangle is in equilibrium when θ = 1 \pi, show that k = 4\sqrt\{3\} + 6.\]
The triangle is released from rest from the position at which θ = 1 \pi. Show that when it passes\par
through the equilibrium position its speed V satisfies\par
4ag\par
\[V 2 = (6 + \sqrt\{3\}) .\]
\end{flushleft}
\subsection*{Page 8}
\begin{flushleft}
\bigskip\textbf{Section C: Probability and Statistics}\par
\bigskip\textbf{Question 12.} A list consists only of letters A and B arranged in a row. In the list, there are a letter As and\par
b letter Bs, where a \leq 2 and b \leq 2, and a + b = n. Each possible ordering of the letters is\par
equally probable. The random variable X is defined by\par
1 if the first letter in the row is A;\par
\[X =\]
\}0 otherwise.\par
\[The random variables X (2 \geq k \geq n) are defined by\]
k\par
1 if the (k 1)th letter is B and the kth is A;\par
\[X = -\]
k\par
\}0 otherwise.\par
n\par
\[The random variable S is defined by S = X .\]
i\par
\[i=1\]
(\par
\textbf{(i)} Find expressions for E(X ), distinguishing between the cases i = 1 and i = 1, and show\par
i\par
\[\geq\]
a(b + 1)\par
\[that E(S) = .\]
n\par
\textbf{(ii)} Show that:\par
a(a 1)b\par
\[(a) for j \leq 3, E(X X ) = - ;\]
1 j\par
n(n 1)(n 2)\par
- -\par
n-2 n a(a 1)b(b 1)\par
\[(b) E(X X ) = - - ;\]
i j\par
2n(n 1)\par
\[i=2 j=i+2\]
\[\int \sum -\]
( (\par
a(a 1)b(b + 1)\par
\[(c) Var (S) = - .\]
\[n^2(n 1)\]
-\par
\end{flushleft}
\subsection*{Page 9}
\begin{flushleft}
13 (a) The continuous random variable X satisfies 0 \geq X \geq 1, and has probability density\par
function f(x) and cumulative distribution function F(x). The greatest value of f(x)\par
\[is M , so that 0 \geq f(x) \geq M .\]
\textbf{(i)} Show that 0 \geq F(x) \geq M x for 0 \geq x \geq 1.\par
\textbf{(ii)} For any function g(x), show that\par
1 1\par
\[\geq 2\]
\[2g(x)F(x)f(x)dx = g(1) g (x) F(x) dx .\]
-\par
\[\leq0 \leq0\]
) [\par
(b) The continuous random variable Y satisfies 0 \geq Y \geq 1, and has probability density\par
function kF(y)f(y), where f and F are as above.\par
\textbf{(i)} Determine the value of the constant k.\par
\textbf{(ii)} Show that\par
nM nM\par
\[1 + μ \geq E(Y n) \geq 2M μ ,\]
n+1 n+1\par
n + 1 - n + 1\par
\[where μ = E(Xn+1) and n \leq 0.\]
n+1\par
\textbf{(iii)} Hence show that, for n \leq 1,\par
n n 1\par
\[μ \leq - .\]
n\par
(n + 1)M - n + 1\par
\end{flushleft}
\end{document}
